\documentclass[11pt]{article}
\usepackage[margin=1in]{geometry}
\usepackage[hidelinks]{hyperref}

\title{Sri Lankan University Course Eligibility Expert System}
\author{}
\date{}

\begin{document}
\maketitle

\section{Scope}

This Python command-line application checks whether supplied applicant facts
satisfy the minimum eligibility rules represented for nine Sri Lankan
state-university courses. It does not predict admission, selection, ranking,
Z-scores, quotas, cut-offs, or available places.

\section{Knowledge Base and Source}

The knowledge base contains rules R01--R21 independently of the command-line
interface. Every rule cites the source identifier
\texttt{UGC-SL-ADMISSIONS-HANDBOOK-2025-2026}, representing the University
Grants Commission Sri Lanka publication \emph{Admission to Undergraduate
Courses of the Universities in Sri Lanka, Academic Year 2025/2026}. The
official publication is available at
\url{https://www.ugc.ac.lk/downloads/admissions/Handbook_2025_26/student_handbook_english.pdf}.

The stored page numbers are the verified printed pages: R01--R03 page 9; R04
page 13; R05 page 50; R06--R07 page 51; R09--R11 page 57; R12--R13 page 58;
R08 page 68; R14 page 83; R18--R21 page 84; and R15--R17 page 90.

\section{Inference Design}

The inference layer provides two explicit workflows. The implementation is a
small evaluator for a fixed rule map; it is not Prolog, RETE, or a fully
general-purpose inference engine.

\subsection{Backward Chaining}

For a single-course request, \texttt{backward\_chaining()} begins with the goal
``eligible for selected course.'' It identifies the general rules R01--R04 and
only the course-specific rules needed for that goal. Each required antecedent
is evaluated against the applicant facts. The returned trace records the goal,
required rules, evaluated rule results, failed conditions, source references,
and final conclusion.

\subsection{Forward Chaining}

For an all-courses request, \texttt{forward\_chaining()} begins with the
applicant facts, evaluates R01--R21 once, and derives an eligibility conclusion
for every supported course. The returned trace records the inference mode,
evaluated rules, derived conclusions, failed conditions, and source references.

\section{Explanation and Limitations}

Every course explanation displays the inference mode, final result, general
and course rule outcomes, all failed requirements, and grouped source/page
references. Results depend on accurate and complete input. Missing facts fail
requirements that cannot be established, subject names must match the
documented names, and handbook rules must be reviewed when the academic year
changes.

\end{document}
